% TOP_PROBLEM: {"id": 30007349, "title": "The nested CLE limit of half-plane double-dimer loops", "category_id": 19, "category": {"id": 19, "name": "probability", "display_name": "Probability", "description": "Problems involving probability theory, stochastic processes, and random structures.", "slug": "probability", "order_index": 19, "created_at": "2026-07-31T15:26:25.671Z"}, "status": "open"}
% FIRST_POSTED: 2026-10-06
% ATTEMPT_STATUS: solved
% OPENAI_CLAIM: {"schema_version":1,"source_repo":"https://github.com/openai/math","source_commit":"adc7f1241b42e322a6451854ab7e4b4c146bf78a","release_date":"2026-10-06","match":"full","resolution":"proved","independently_reviewed":false,"second_pass":"agreed","families":[{"family":"226","title":"The double-dimer loop ensemble converges to CLE<sub>4</sub>","subject":"Probability and statistical mechanics","manuscripts":[{"title":"The Curve Scaling Limit of Half-Plane Double Dimers","pdf_path":"preprints/The-curve-scaling-limit-of-half-plane-double-dimers-September-23-2026/paper.pdf","readme_path":"preprints/The-curve-scaling-limit-of-half-plane-double-dimers-September-23-2026/README.md","date":"2026-09-23","theorem_ref":"thm:main (sections/10-introduction.tex:84)","statement_sources":[{"path":"preprints/The-curve-scaling-limit-of-half-plane-double-dimers-September-23-2026/build/sections/10-introduction.tex","line":84,"label":"thm:main"}]}],"lean":null,"reasoning_trace":null}]}
\documentclass[11pt]{article}
\usepackage[margin=25mm]{geometry}
\usepackage{fontspec}
\setmainfont{Latin Modern Roman}
\usepackage{amsmath,amssymb,mathtools}
\usepackage{unicode-math}
\setmathfont{Latin Modern Math}
\usepackage{enumitem}
\usepackage{xurl,hyperref,bookmark}
\hypersetup{colorlinks=true,urlcolor=blue}
\setcounter{secnumdepth}{0}
\setlength{\parindent}{0pt}
\setlength{\parskip}{5pt}
\setlength{\emergencystretch}{3em}
\begin{document}
\section*{The nested CLE limit of half-plane double-dimer loops}

\subsection{OpenAI's claim}
OpenAI presents a proof in the selected manuscript. All macroscopic half-plane Temperleyan double-dimer loops are claimed to converge to the nested $\mathrm{CLE}_4$ ensemble along the full mesh limit.

\subsection{Mathematical statement of the claimed result}
\paragraph{Result 1.} On the upper-half-plane square lattice, use the Temperleyan matching law obtained by exhausting with coarse-grid simple polygons, deleting a boundary root escaping to infinity, and taking local limits of uniform perfect matchings. Superpose two independent samples and remove doubled edges, leaving loops \(\mathcal L_\delta\). For \(F(z)=(z-i)/(z+i)\), the source claims tightness and
\[
 F(\mathcal L_\delta)\Longrightarrow F(\mathcal L),
\]
where \(\mathcal L\) is nested \(\operatorname{CLE}_4\) in the upper half-plane, normalized by Brownian-loop-soup central charge one. Convergence matches all loops above every positive diameter threshold in uniform curve distance modulo circle reparametrizations, retaining multiplicity and nesting, along the full mesh limit.

\subsection{Outline of the argument}
Rectangle transfer is expressed in particle sectors and controlled by singular modes with marked boundary caps. Traversal and short-segment bounds provide curve tightness; the resulting geometric control upgrades the dimer field information to matching of every macroscopic loop.

\subsection{Scope relative to this problem}
All macroscopic half-plane Temperleyan double-dimer loops are claimed to converge to the nested $\mathrm{CLE}_4$ ensemble along the full mesh limit. The question retains these same parameters and conventions. Within these assumptions, this is a claimed affirmative answer to the stated question.

\subsection{Formal verification}
Family 226 has no Lean scope page in this release. Its selected result is not presented as formalized. This site has not executed the Lean code. The scope descriptions do not certify the remaining manuscript arguments.

\subsection{Status caveat}
OpenAI attributes these manuscripts to work by an internal model. Its release cautions that arguments without formalization can contain mistakes. This catalogue has not independently checked the proofs; the labels report the mathematical claims attributed to OpenAI.

\subsection{Sources}
\begin{itemize}
\item OpenAI, \emph{The Curve Scaling Limit of Half-Plane Double Dimers}, 2026-09-23; citation key \texttt{OAI:The-curve-scaling-limit-of-half-plane-double-dimers-September-23-2026}. \href{https://github.com/openai/math/blob/adc7f1241b42e322a6451854ab7e4b4c146bf78a/preprints/The-curve-scaling-limit-of-half-plane-double-dimers-September-23-2026/paper.pdf}{PDF (GitHub)}; \href{https://github.com/openai/math/raw/adc7f1241b42e322a6451854ab7e4b4c146bf78a/preprints/The-curve-scaling-limit-of-half-plane-double-dimers-September-23-2026/paper.pdf}{PDF (direct)}; \href{https://github.com/openai/math/blob/adc7f1241b42e322a6451854ab7e4b4c146bf78a/preprints/The-curve-scaling-limit-of-half-plane-double-dimers-September-23-2026/README.md}{Manuscript page and citation}; \href{https://github.com/openai/math/blob/main/preprints/The-curve-scaling-limit-of-half-plane-double-dimers-September-23-2026/paper.pdf}{Latest version}; Statement source: \href{https://github.com/openai/math/blob/adc7f1241b42e322a6451854ab7e4b4c146bf78a/preprints/The-curve-scaling-limit-of-half-plane-double-dimers-September-23-2026/build/sections/10-introduction.tex\#L84}{thm:main}.
\item \href{https://github.com/openai/math/blob/adc7f1241b42e322a6451854ab7e4b4c146bf78a/CONTENTS.md}{OpenAI catalogue entry}.
\end{itemize}
\end{document}
